% 6. ALP の自己相互作用
\section{ALP の自己相互作用}\label{sec:selfint}

\subsection{ALP のポテンシャル}

ALP(axion-like particle)は、周期的なポテンシャルをもつスカラー場である。
\begin{equation}
V(\varphi)=\mu^2f^2\Bigl[1-\cos\Bigl(\frac{\varphi}{f}\Bigr)\Bigr].
\label{eq:Vcos}
\end{equation}
$f$ を崩壊定数と呼ぶ。$\cos x=1-\dfrac{x^2}{2}+\dfrac{x^4}{24}-\cdots$ を $x=\varphi/f$ で使うと
\begin{equation}
V(\varphi)=\mu^2f^2\Bigl[\frac{\varphi^2}{2f^2}-\frac{\varphi^4}{24f^4}+\cdots\Bigr]
=\frac12\mu^2\varphi^2-\frac{\lambda}{4!}\varphi^4+\cdots,
\qquad \lambda\equiv\frac{\mu^2}{f^2}.
\label{eq:Vexpand}
\end{equation}
第 1 項が質量の項、第 2 項が四次の自己相互作用である。符号が負なので、粒子どうしは引き合う。$f$ が小さいほど自己相互作用は強い。

\paragraph{QCD アクシオンと ALP}
QCD アクシオンでは $\mu$ と $f$ が $\mu\approx0.57\,\mathrm{neV}\times(10^{16}\GeV/f)$ で結ばれていて、$\mu=0.71\times10^{-12}\eV$ なら $f\approx8\times10^{18}\GeV$ になる。ALP では $\mu$ と $f$ は独立で、研究の窓は $f\sim10^{12}\GeV$ にある。

\paragraph{$\lambda$ の大きさ}
$\mu=7\times10^{-13}\eV$、$f=10^{12}\GeV=10^{21}\eV$ では $\lambda=\mu^2/f^2=(7\times10^{-34})^2=4.9\times10^{-67}$。極端に小さいが、雲の粒子数は $N\sim10^{77}$ なので、まとまった効果が大きくなる。

\subsection{運動方程式と非線形 Schr\"odinger 方程式}

運動方程式は $\partial_t^2\varphi-\nabla^2\varphi+V'(\varphi)=0$ で、$V'=\mu^2\varphi-(\lambda/6)\varphi^3$ だから
\begin{equation}
\partial_t^2\varphi-\nabla^2\varphi+\mu^2\varphi-\frac{\lambda}{6}\varphi^3=0.
\label{eq:EOM_lambda}
\end{equation}
重力を入れるのは \S\ref{sec:atom} と同じなので、ここでは新しい $\varphi^3$ の項だけを扱う。

\paragraph{$\varphi^3$ の展開}
$a\equiv\psi e^{-i\mu t}$ と書くと、$\varphi=(2\mu)^{-1/2}(a+a^*)$ で
\begin{equation}
\varphi^3=(2\mu)^{-3/2}\bigl(a^3+3a^2a^*+3aa^{*2}+a^{*3}\bigr).
\end{equation}
$a^2a^*=|\psi|^2\psi\,e^{-i\mu t}$ は、もとの場と同じ $e^{-i\mu t}$ で振動する。$a^3\propto e^{-3i\mu t}$ は速く振動し、ゆっくり変わる $\psi$ の式には効かないので落とす。$e^{-i\mu t}$ の部分は
\begin{equation}
\varphi^3\supset(2\mu)^{-3/2}\cdot3|\psi|^2\psi\,e^{-i\mu t}.
\end{equation}

\paragraph{集める}
\S\ref{sec:atom} と同じ手順で $e^{-i\mu t}$ の係数を集めると(重力の項を含めて)
\begin{equation}
(2\mu)^{-1/2}\bigl(-2i\mu\dot\psi-\nabla^2\psi+2\Phi\mu^2\psi\bigr)-\frac{\lambda}{6}(2\mu)^{-3/2}\cdot3|\psi|^2\psi=0.
\end{equation}
両辺に $(2\mu)^{1/2}$ をかけると
\begin{equation}
-2i\mu\dot\psi-\nabla^2\psi+2\Phi\mu^2\psi-\frac{\lambda}{2}\cdot\frac{|\psi|^2\psi}{2\mu}=0.
\end{equation}
$-2\mu$ で割って整理する($\mu\Phi=-\alpha/r$)。
\begin{equation}
\boxed{\;i\frac{\partial\psi}{\partial t}=\Bigl(-\frac{\nabla^2}{2\mu}-\frac{\alpha}{r}\Bigr)\psi-\frac{\lambda}{8\mu^2}|\psi|^2\psi\;}
\label{eq:GP}
\end{equation}
これは引力の非線形 Schr\"odinger 方程式(Gross--Pitaevskii 方程式)で、\arxiv{2011.11646}(以下 BGLS)の式 12 と一致する。

\subsection{非線形項は準位をつなぐ}

雲がいくつかの準位の重ね合わせだとする。
\begin{equation}
\psi=\sum_i c_i\,\Psi_i(\bm x)\,e^{-i(E_i-\mu)t},\qquad \Psi_i\propto e^{im_i\phi}.
\end{equation}
非線形項 $|\psi|^2\psi=\psi\,\psi\,\psi^*$ には、3 つの準位の積 $c_ic_jc_k^*\Psi_i\Psi_j\Psi_k^*$ が現れる。この積は
\begin{itemize}
\item 振動数 $E_i+E_j-E_k$ で振動し、
\item 方位の依存性 $e^{i(m_i+m_j-m_k)\phi}$ をもつ。
\end{itemize}
これは、粒子 $i$ と $j$ が衝突して $k$ ともう一つの状態に移る過程にあたる。もう一つの状態は、振動数 $\omega'=E_i+E_j-E_k$、方位量子数 $m'=m_i+m_j-m_k$ をもつ。その行き先は振動数で決まる。
\begin{itemize}
\item $\omega'>\mu$: 束縛されず、無限遠に逃げる(スカラー波の放出)。
\item $\omega'<\mu$ で、育つ準位に一致しない: 束縛されているが超放射しないモードを強制的に揺らし、そのモードは BH に吸収される。
\end{itemize}

\paragraph{主な二つの過程}
準位のエネルギー $E_n=\mu(1-\alpha^2/2n^2)$ を使って計算する。

\textbf{(a) $211\times211\to322\times\mathrm{BH}$}: $i=j=211$、$k=322$。
\begin{align}
\omega'&=2E_{211}-E_{322}=\mu\Bigl[2\Bigl(1-\frac{\alpha^2}{8}\Bigr)-\Bigl(1-\frac{\alpha^2}{18}\Bigr)\Bigr]
=\mu\Bigl(1-\frac{\alpha^2}{4}+\frac{\alpha^2}{18}\Bigr)=\mu\Bigl(1-\frac{7\alpha^2}{36}\Bigr)<\mu,\\
m'&=1+1-2=0.
\end{align}
束縛されていて $m'=0$ なので超放射せず、BH に吸収される。結果として、$|211\rangle$ の粒子 2 個が、$|322\rangle$ の 1 個と、BH に落ちる 1 個になる。BH に落ちる粒子は角運動量をもたない($m'=0$)。

\textbf{(b) $322\times322\to211\times\infty$}: $i=j=322$、$k=211$。
\begin{align}
\omega'&=2E_{322}-E_{211}=\mu\Bigl[2\Bigl(1-\frac{\alpha^2}{18}\Bigr)-\Bigl(1-\frac{\alpha^2}{8}\Bigr)\Bigr]
=\mu\Bigl(1-\frac{\alpha^2}{9}+\frac{\alpha^2}{8}\Bigr)=\mu\Bigl(1+\frac{\alpha^2}{72}\Bigr)>\mu,\\
m'&=2+2-1=3.
\end{align}
無限遠に逃げる。$|322\rangle$ の粒子 2 個が、$|211\rangle$ の 1 個と、角運動量 $3\hbar$ をもって逃げる 1 個になる。

\begin{figure}[H]
\centering
% 自己相互作用による二つの過程(模式図)
\begin{tikzpicture}[
  st/.style={draw, rounded corners=2pt, minimum width=1.2cm, minimum height=0.6cm, font=\small, line width=0.7pt},
  ar/.style={-{Stealth[length=2.2mm]}, line width=0.8pt}]
  % (a)
  \node[font=\small\bfseries] at (2.0,2.3) {(a) $211\times211\to322\times\mathrm{BH}$};
  \node[st, draw=cA, text=cA] (a1) at (0,1.2) {211};
  \node[st, draw=cA, text=cA] (a2) at (0,0.0) {211};
  \node[circle, fill=ink2, inner sep=1.5pt] (v) at (1.6,0.6) {};
  \node[st, draw=cB, text=cB] (a3) at (3.4,1.2) {322};
  \node[st, draw=ink, fill=ink, text=white] (bh) at (3.4,0.0) {BH};
  \draw[ar, cA] (a1) -- (v); \draw[ar, cA] (a2) -- (v);
  \draw[ar, cB] (v) -- (a3);
  \draw[ar, ink, decorate, decoration={snake, amplitude=1.2pt, segment length=5pt, post length=2mm}] (v) -- (bh);
  \node[font=\scriptsize, color=ink2, anchor=north] at (2.5,-0.45) {$\omega'=\mu(1-7\alpha^2/36)<\mu$、$m'=0$};
  % (b)
  \begin{scope}[xshift=7.2cm]
  \node[font=\small\bfseries] at (2.0,2.3) {(b) $322\times322\to211\times\infty$};
  \node[st, draw=cB, text=cB] (b1) at (0,1.2) {322};
  \node[st, draw=cB, text=cB] (b2) at (0,0.0) {322};
  \node[circle, fill=ink2, inner sep=1.5pt] (w) at (1.6,0.6) {};
  \node[st, draw=cA, text=cA] (b3) at (3.4,1.2) {211};
  \node[st, draw=cC, text=cC!70!black] (inf) at (3.4,0.0) {$\infty$};
  \draw[ar, cB] (b1) -- (w); \draw[ar, cB] (b2) -- (w);
  \draw[ar, cA] (w) -- (b3);
  \draw[ar, cC!70!black, decorate, decoration={snake, amplitude=1.2pt, segment length=5pt, post length=2mm}] (w) -- (inf);
  \node[font=\scriptsize, color=ink2, anchor=north] at (2.5,-0.45) {$\omega'=\mu(1+\alpha^2/72)>\mu$、$m'=3$};
  \end{scope}
\end{tikzpicture}

\caption{自己相互作用による二つの過程。(a) $|211\rangle$ の 2 個が $|322\rangle$ と BH に吸収される $m=0$ の波になる。(b) $|322\rangle$ の 2 個が $|211\rangle$ と無限遠に逃げる $m=3$ の波になる。}
\label{fig:processes}
\end{figure}

\subsection{Bose 増強と率の方程式}

\paragraph{Fermi の黄金律と占有数}
ボソンの遷移 $i+j\to k+l$ の率は、Fermi の黄金律で $|\text{行列要素}|^2$ に比例し、さらに始状態と終状態の占有数の因子がつく。
\begin{equation}
\text{率}\propto N_i\,N_j\,(1+N_k)\,(1+N_l).
\end{equation}
$(1+N)$ は誘導放出の因子(Bose 増強)である。過程 (a) では $l$ は BH に落ちる状態で $N_l\simeq0$、$k=322$ は占有数が巨大なので $(1+N_{322})\simeq N_{322}$。したがって
\begin{equation}
\text{(a) の率}=\Gamma_\mathrm{BH}\,N_{211}^2\,N_{322},\qquad
\text{(b) の率}=\Gamma_\infty\,N_{322}^2\,N_{211}.
\end{equation}
$N\sim10^{77}$ の雲では古典的な波の非線形性と同じだが、粒子の言葉では「誘導された 3 体の遷移」と読める。

\paragraph{率の方程式}
1 回の (a) で $N_{211}$ は 2 減り $N_{322}$ は 1 増える。1 回の (b) で $N_{322}$ は 2 減り $N_{211}$ は 1 増える。超放射($\Gamma_{211}$、$\Gamma_{322}$、粒子数の率)と合わせて
\begin{align}
\frac{\dd N_{211}}{\dd t}&=\Gamma_{211}N_{211}-2\Gamma_\mathrm{BH}N_{211}^2N_{322}+\Gamma_\infty N_{211}N_{322}^2,
\label{eq:rate211}\\
\frac{\dd N_{322}}{\dd t}&=\Gamma_{322}N_{322}+\Gamma_\mathrm{BH}N_{211}^2N_{322}-2\Gamma_\infty N_{211}N_{322}^2,
\label{eq:rate322}\\
\frac{\dd J}{\dd t}&=-\hbar\bigl(\Gamma_{211}N_{211}+2\Gamma_{322}N_{322}\bigr).
\label{eq:rateJ}
\end{align}
BH の角運動量を変えるのは超放射だけである。(a) で BH に落ちる波は $m=0$ なので角運動量を運ばない。これは BGLS の式 34--36 の主な項である(重力波の放出の項は小さいので省いた)。

\subsection{率の係数と、$\varepsilon$ の単位}

$\Gamma_\mathrm{BH}$、$\Gamma_\infty$ は行列要素の 2 乗なので $\lambda^2$ に比例し、残りは $\alpha$ のべきと数値の係数になる。BGLS は水素原子の波動関数で重なり積分を計算した(式 18、21、Table I)。
\begin{equation}
\Gamma_\mathrm{BH}\simeq4\times10^{-7}\,\alpha^7\lambda^2\rp\,\mu,\qquad
\Gamma_\infty\simeq1\times10^{-8}\,\alpha^4\lambda^2\,\mu.
\label{eq:GammaN_SI}
\end{equation}

\paragraph{$\varepsilon$ の単位に直す}
式 \eqref{eq:epsilon} の $\varepsilon=N/(GM^2)$($\hbar=c=1$)を使う。$N=GM^2\,\varepsilon$ を式 \eqref{eq:rate211} に代入して $GM^2$ で割ると、3 体の項の係数は $(GM^2)^2$ 倍になる。
\begin{equation}
\gamma_\mathrm{BH}\equiv\Gamma_\mathrm{BH}\,(GM^2)^2,\qquad \gamma_\infty\equiv\Gamma_\infty\,(GM^2)^2.
\end{equation}
$\lambda=\mu^2/f^2$、$G=1/\Mpl^2$、$\alpha=GM\mu=M\mu/\Mpl^2$ を使うと
\begin{equation}
\lambda^2(GM^2)^2=\frac{\mu^4}{f^4}\cdot\frac{M^4}{\Mpl^4}=\Bigl(\frac{M\mu}{\Mpl^2}\Bigr)^4\Bigl(\frac{\Mpl}{f}\Bigr)^4=\alpha^4\Bigl(\frac{\Mpl}{f}\Bigr)^4.
\end{equation}
したがって
\begin{equation}
\boxed{\;\gamma_\mathrm{BH}=\kappa_\mathrm{BH}\,\alpha^{11}\Bigl(\frac{\Mpl}{f}\Bigr)^4\rp\,\mu,\qquad
\gamma_\infty=\kappa_\infty\,\alpha^{8}\Bigl(\frac{\Mpl}{f}\Bigr)^4\mu\;}
\label{eq:gamma_SI}
\end{equation}
で、BGLS の精密な値(Table V・VI)は $\kappa_\mathrm{BH}=4.3\times10^{-7}$、$\kappa_\infty=1.1\times10^{-8}$。$\Mpl=1/\sqrt G=1.22\times10^{19}\GeV$ である。$\varepsilon$ の単位での率の方程式は、式 \eqref{eq:rate211}--\eqref{eq:rateJ} の $N$ を $\varepsilon$ に、$\Gamma$ を $\gamma$ に置き換えたもので、角運動量の式は $\dd\chi/\dd t=-\gamma_{211}\varepsilon_{211}-2\gamma_{322}\varepsilon_{322}$ になる($\chi=J/GM^2$、質量を固定)。超放射の率は単位によらず $\gamma_{211}=\Gamma_{211}$ である。

\begin{caution}
\begin{itemize}
\item 式 \eqref{eq:gamma_SI} は $\alpha$ の主要項である。Kerr 時空での計算(\arxiv{2412.03655})では、$\gamma_\mathrm{BH}$ は $\alpha=0.1$ で約 2.5 倍、$\alpha=0.2$ で約 5 倍大きい(研究ノート E7)。$\gamma_\infty$ の補正は 1.0--1.2 倍で小さい。
\item $\Mpl$ を換算 Planck 質量 $2.4\times10^{18}\GeV$ と取り違えると、$(\Mpl/f)^4$ が 35 倍ずれる。
\end{itemize}
\end{caution}

\begin{summary}
\begin{itemize}
\item ALP のポテンシャルを展開すると引力の四次の結合 $-(\lambda/4!)\varphi^4$、$\lambda=\mu^2/f^2$ が出る。非相対論の極限では $-(\lambda/8\mu^2)|\psi|^2\psi$ の項になる。
\item 非線形項は準位をつなぐ。$211\times211\to322\times\mathrm{BH}$(振動数 $\mu(1-7\alpha^2/36)$、$m=0$、BH に吸収)と $322\times322\to211\times\infty$(振動数 $\mu(1+\alpha^2/72)$、$m=3$、無限遠へ)。
\item 率は Bose 増強で $N_{211}^2N_{322}$、$N_{322}^2N_{211}$ に比例する。$\varepsilon$ の単位で $\gamma_\mathrm{BH}\propto\alpha^{11}(\Mpl/f)^4$、$\gamma_\infty\propto\alpha^8(\Mpl/f)^4$。
\end{itemize}
\end{summary}

\begin{tocode}
コード: \texttt{si.g\_bh}、\texttt{si.g\_inf} が式 \eqref{eq:gamma_SI}(\texttt{KAPPA\_BH}、\texttt{KAPPA\_INF}、\texttt{MPL\_GEV})。\texttt{e05\_kinetics.py} が率の方程式を ODE で解く(研究ノート E5)。
\end{tocode}
